% SPDX-License-Identifier: Apache-2.0
% covers: AxiPerPins, AxiPerDriven
\datasheet{AxiPerPins}{The AXI link's peripheral end joined to an AXI4 peripheral's pins.}

\dsfacts{\code{lib/parts/src/bus/axi\_per\_pins.rs}}%
{\code{txhdl\_parts::bus::axi\_per\_pins::AxiPerPins}}%
{\code{//docs:pcie}}

\subsection*{Function}

A core from elsewhere that is an AXI4 peripheral, such as AMD's DDR3
controller generated with its AXI4 port, has pins: a \code{valid} and a
\code{ready} per channel and a wire per field. \code{AxiPerPins} is
\code{AxiPins} the other way round. It stands on the link's peripheral
end, where an \code{AxiPer} would, and puts the link's channels onto
those pins, holding nothing.

Its inputs are the side \code{inp: AxiPerPinsIn}: the eleven pins the
peripheral drives as one field, \code{pins}, an \code{AxiPerDriven},
and the \code{aw}, \code{ar} and \code{w} channels from the link, so
its input ports are \code{inp\_pins\_awready} and the rest beside
\code{inp\_aw}, \code{inp\_ar} and \code{inp\_w}. \code{AxiPerDriven}
derives \code{Ports}, so a design holding a peripheral's pins among its
own ports passes them on whole. Its outputs are the side
\code{outp: AxiPerPinsOut}: the \code{b} and \code{r} channels back to
the link, and the twenty-six pins the peripheral reads.

\code{axi\_per\_pins::sim} has \code{PinRam}, a memory on the pins
written as a simulation, and \code{pins}, which makes the wires.
\code{PinRam} can take a warm-up, before which it is not ready, and a
read and a write latency, and it can take an address past its end
modulo its size; that is how \code{ddr3::Ddr3} models AMD's controller.

\subsection*{Parameters}

\begin{center}\footnotesize
\begin{tabular}{@{}lp{0.7\textwidth}@{}}
\toprule
Parameter & Meaning \\
\midrule
\code{A} & The address width, in bits. The tables use 32. \\
\code{D} & The data width, in bits. The tables use 32. \\
\code{S} & The strobe width, \code{D / 8}. The tables use 4. \\
\code{I} & The identifier width, in bits. The tables use 5. \\
\bottomrule
\end{tabular}
\end{center}

The tables are generated for \code{AxiPerPins<32, 32, 4, 5>}, the
widths \code{Ddr3Per} uses it with on the board. The example
\code{ex\_axi\_per\_pins} uses \code{AxiPerPins<30, 32, 4, 5>}, the
controller's own 30-bit address.

\dstables{AxiPerPins}

\subsection*{Behaviour}

\begin{itemize}
\item It holds no state; every output is a function of the step's
inputs and the channels.
\item An address phase at the head of \code{aw} is \code{awvalid} and
the \code{aw} pins: \code{awid}, \code{awaddr}, \code{awlen},
\code{awsize}, the burst type encoded in two bits, \code{awlock},
\code{awcache}, \code{awprot} and \code{awqos}. It is taken off the
channel in a step \code{awready} is high. The same holds for
\code{ar}, and for \code{w} with \code{wdata}, \code{wstrb} and
\code{wlast}.
\item A response the peripheral offers, \code{bvalid} with
\code{bid} and \code{bresp}, the response decoded from its two bits,
goes onto \code{b} in a step \code{b} has room, and \code{bready} is
that room. The same holds for \code{rvalid} and \code{r}, with
\code{rid}, \code{rdata}, \code{rresp} and \code{rlast}, and
\code{rready}.
\item Every field of the link's address phase reaches a pin except
\code{region}: AMD's controller has no region pins, and a peripheral
joined to one link needs none, since the region is how an interconnect
tells decoded ranges of one peripheral apart. \code{AxQOS} is forwarded.
\end{itemize}

\subsection*{Verification}

\begin{itemize}
\item \code{a\_host\_on\_the\_link\_reaches\_a\_peripheral\_on\_pins} in
\code{bus::axi\_per\_pins}, in \code{//lib/parts:parts\_test}, writes
three words through an \code{AxiHost}, the unit and a \code{PinRam},
reads four back across the burst's ends, and reads past the memory's
end, answered \code{SlvErr}.
\item \code{a\_timed\_memory\_waits\_and\_then\_answers\_late} there
holds \code{PinRam}'s timing: a warm-up holds a write off, and a latency
delays an answer by itself and no more.
\item \code{ex\_axi\_per\_pins} runs \code{AxiPerPins<30, 32, 4, 5>} with
a host on the link and a \code{PinRam} on the pins, and the build
simulates its netlist against that run under nvc and Verilator:
\code{//docs:axi\_per\_pins\_sim\_axi\_per\_pins\_tb\_test} and
\code{//docs:axi\_per\_pins\_vsim\_test}.
\item \code{//ddr3:ddr3\_test} drives it inside \code{Ddr3Per}, and
\code{//ddr3/sim:wrapper\_test}, which is manual, checks the
controller's side of the same pins under Vivado's simulator.
\end{itemize}

\subsection*{Used by}

\code{Ddr3Per}, as its field \code{pins}, between the board's router
and AMD's DDR3 controller. The example \code{ex\_axi\_per\_pins}.

\subsection*{Limits}

Its \code{ready} outputs are the channels' room in the same step, so
a path runs from a channel's state through the unit to the
peripheral's pins with no register. It has no \code{AxREGION} or user
pins, so a peripheral that wants a region gets it as a constant where
it is instantiated.
